% ID: SEL-D0020-OB-E03
% ESTADO: cotejada. Segunda pasada visual completa y borrador compilado revisados.
% Fuente: Selectividad/matematicas_aplicadas/reserva b/Examen-Matemáticas Aplicadas a las Ciencias Sociales.pdf
% SHA-256: 37584a9604930b74008a6bda88ef3d61fa926d5e6a540ab9013569258f772256
% Página: 2; opción B; ejercicio original 3.

Una almazara recibe cajas de aceitunas de dos productoras, A y B, que cultivan dos variedades, picual y arbequina. El 40\% proviene de la productora A, de las cuales el 60\% es de la variedad picual. De las que provienen de la productora B, el 30\% es de la variedad arbequina. Se elige una caja de aceitunas al azar.
\begin{enumerate}[label=\alph*)]
\item \textbf{(1punto)} \textquestiondown{}Cuál es la probabilidad de que sea de la variedad picual?
\item \textbf{(1 punto)} Si se sabe que es de la variedad picual, \textquestiondown{}cuál es la probabilidad de que provenga de la productora A?
\item \textbf{(0.5 puntos)} Calcule la probabilidad de que sea de la productora A o de la variedad picual.
\end{enumerate}
